\documentclass[10pt,a4paper]{article}

% Mise en page et typographie
\usepackage[T1]{fontenc}
\usepackage[utf8]{inputenc}
\usepackage[french]{babel}
\usepackage{lmodern}
\usepackage{microtype}
\usepackage[a4paper,margin=2cm]{geometry}
\usepackage{array}
\usepackage{longtable}
\usepackage{booktabs}
\usepackage{pdflscape}
\usepackage{enumitem}
\usepackage{hyperref}
\usepackage{fancyhdr}
\usepackage{lastpage}

\hypersetup{
  hidelinks,
  pdftitle={Psychophysique de la détection des perturbations internes dans les LLM},
  pdfauthor={William Auroux, Thomas Gegout, Louis Le Dain et Seif Zaafouri}
}

\setlength{\parindent}{0pt}
\setlength{\parskip}{5pt}
\setlist[itemize]{leftmargin=1.4em,itemsep=1pt,topsep=2pt}

\pagestyle{fancy}
\fancyhf{}
\renewcommand{\headrulewidth}{0pt}
\renewcommand{\footrulewidth}{0pt}
\fancyfoot[C]{\footnotesize \thepage\,/\,\pageref{LastPage}}

\newcolumntype{L}[1]{>{\raggedright\arraybackslash}p{#1}}

\begin{document}

\begin{center}
  \vspace*{8pt}
  {\LARGE\bfseries Psychophysique de la détection des perturbations internes dans les LLM :\\[-1pt]
  Introspection, anomalie ou géométrie ?\par}
  \vspace{10pt}
  {\normalsize\bfseries William AUROUX \hspace{1.3em} Thomas GEGOUT \hspace{1.3em} Louis LE DAIN \hspace{1.3em} Seif ZAAFOURI\par}
  \vspace{4pt}
  {\small\itshape Encadrant : Fabrice POPINEAU\par}
  \vspace{3pt}
  {\small 7 septembre 2026\par}
\end{center}

\vspace{5pt}
\section*{Synthèse bibliographique}

Avec la croissance rapide des capacités des grands modèles de langage au cours des dernières années, les questions de sécurité et d'interprétabilité de l'intelligence artificielle (IA) sont d'autant plus pertinentes. Être capable de comprendre comment un modèle se comporte et pourquoi représente un domaine de recherche majeur. En particulier, la question de l'introspection des LLM est intéressante pour améliorer la fiabilité et l'alignement des systèmes d'IA, et permet au sens large d'interroger l'attribution de capacités cognitives et psychiques humaines aux modèles. La présence d'un tel phénomène permettrait à terme de clarifier plus simplement les intentions des modèles de langage pour des utilisateurs humains.

La méthode d'évaluation de l'introspection des LLM par injection de vecteurs de concept, ou \emph{steering}, initialement conçue par Lindsey \cite{lindsey2026} selon des critères stricts (exactitude, ancrage causal, internalité et représentation de second ordre), a depuis été retravaillée par Hahami et al. \cite{hahami2026} pour proposer un cadre de mesure plus rigoureux et robuste aux artefacts, notamment le biais global des logits, via des tâches de discrimination forcée (2AFC) plutôt que de détection binaire. Au-delà de son rôle dans l'étude de l'introspection, cette même technique de \emph{steering} soulève des enjeux pratiques pour la sûreté des IA et la performance en utilisation normale : Mishra et al. \cite{mishra2026} montrent que les états internes induits par \emph{steering} sont presque toujours inaccessibles par un prompt textuel réel, ce qui découple structurellement les vulnérabilités démontrées en accès complet (\emph{white-box}) de celles exploitables en accès boîte noire (\emph{black-box}). Kowalski et al. \cite{kowalski2026} montrent par ailleurs que les modèles peuvent altérer leurs activations internes pour échapper à certains contrôles et ainsi cacher un comportement malveillant. Enfin, Nguyen et al. \cite{nguyen2026} analysent la propension d'une technique de \emph{steering} à altérer le comportement d'un modèle, en améliorant ou en baissant ses performances sur des tâches usuelles.

L'état de l'art n'évalue pas simplement les capacités d'introspection des modèles, mais montre également qu'elles sont améliorables. Macar et al. \cite{macar2026} montrent notamment que ces capacités sont déjà présentes mais peu stimulées dans les modèles de base : le \emph{refusal ablation}, ou \emph{abliteration}, augmente la détection des concepts injectés de 53\,\%, tandis qu'un \emph{trained bias vector} l'améliore de 75\,\% sur des concepts non vus, sans augmentation significative des faux positifs. Macar et al. \cite{macar2026}, ainsi que Kowalski et al. \cite{kowalski2026}, montrent également que cette capacité apparaît ou se renforce à la suite d'un post-entraînement général d'un LLM, avec l'algorithme DPO notamment. Par ailleurs, Ferrara \cite{ferrara2026} montre qu'en partant d'un modèle avec un AUROC de 0,5, donc pratiquement aléatoire, celui-ci peut être explicitement entraîné à détecter si une opération de \emph{steering} est présente ou non. Après cet entraînement, le modèle généralise parfaitement sur 100 nouvelles directions. Comme dans les articles précédents, ce fine-tuning ne crée pas ce signal, déjà présent, mais le rend plus détectable dans les sorties du modèle, en renforçant le chemin de restitution verbale.

L'architecture des modèles influence également les conclusions sur la détection des perturbations conceptuelles et leur localisation au niveau des couches. Lindsey \cite{lindsey2026} affirme que des modèles tels que Opus 4.1 détectent mieux les injections sur des couches profondes, alors que Hahami et al. \cite{hahami2026} montrent que Llama 3.1 8B est plus sensible aux injections en début d'architecture. Cela souligne l'importance de conduire des expériences sur différentes architectures de modèle, afin de pouvoir potentiellement attribuer les phénomènes observés à des divergences architecturales. Kowalski et al. \cite{kowalski2026} renforcent ce constat : à travers leur benchmark et différentes métriques, leur étude montre que les modèles diffèrent fortement dans leur capacité à contrôler leurs activations. Ce contrôle est principalement effectif dans la moitié profonde du réseau et n'est pas forcément lié au nombre de paramètres, tandis que le ciblage d'une couche précise échoue pratiquement systématiquement.

Outre les propriétés intrinsèques du modèle, telles que son entraînement et son architecture, l'étude de l'amplitude est importante pour évaluer l'impact du \emph{steering}. Nguyen et al. \cite{nguyen2026} étudient cette question dans le contexte du dommage collatéral causé par l'activation steering. Leur analyse montre que deux déplacements de même norme peuvent avoir des effets très différents selon leur direction dans l'espace d'activation, ce qui invite à tenir compte de la géométrie de cet espace lors de leur calibration. Bien que leur travail porte sur la préservation des performances plutôt que sur la détection des perturbations internes, cette idée peut être transposée dans le contexte de l'étude des capacités introspectives. En effet, un coefficient brut d'intensité ne suffit pas nécessairement à rendre comparables des injections appliquées dans des directions dont la variabilité naturelle diffère.

Fornasiere et al. \cite{fornasiere2026} soulèvent aussi la question de l'utilisation d'autres perturbations qu'une perturbation conceptuelle. Ils montrent que plusieurs modèles de langage sont capables de détecter l'ajout de bruits gaussiens ou de \emph{dropout}, malgré une performance hétérogène des modèles sur la classification de ces deux perturbations.

Afin de pouvoir attribuer une détection de perturbations à une réelle introspection de la part d'un modèle, Singh et al. \cite{singh2026} énoncent deux conditions pour concevoir des expériences robustes. La tâche considérée ne doit pas pouvoir être résolue grâce à des indices présents dans les entrées (\emph{privileged access}). De plus, le modèle doit construire une méta-représentation, c'est-à-dire une représentation de son propre état interne, pour détecter que ses activations ont été perturbées (\emph{second-order computation}).

Ainsi, les deux conditions décrites par Singh et al. \cite{singh2026} remettent en question les expériences réalisées dans l'état de l'art, dont les résultats sont susceptibles de provenir de calculs simples internes au LLM dépendant des entrées plutôt que d'une réelle capacité d'introspection au sens fort du terme. Cette remise en question est aussi présente chez Mishra et al. \cite{mishra2026}, qui montrent que l'activation steering déplace le flux résiduel hors de l'espace des états atteignables par un prompt textuel : aucun prompt ne peut donc reproduire le comportement induit par l'injection de vecteur. Ce résultat appelle à la prudence quant à l'interprétation des succès du \emph{steering} comme preuve d'interprétabilité ou de mise en évidence d'un phénomène d'introspection.

Pris ensemble, ces travaux montrent que la détection d'une perturbation interne dépend de sa nature, de son intensité, de la couche d'injection, de l'architecture et de l'entraînement du modèle. Ces facteurs ont toutefois été étudiés séparément, à l'aide de protocoles et d'échelles d'intensité différents, ce qui empêche de déterminer si les résultats observés sur les expériences d'introspection proviennent du contenu conceptuel des perturbations, de leur caractère inhabituel ou simplement de leur amplitude relative à la géométrie des activations. Notre travail s'inscrit dans le prolongement de cet état de l'art en réunissant ces dimensions dans un même cadre expérimental contrôlé.

Dans nos travaux, nous proposons d'étudier de manière approfondie différents types de perturbations d'un LLM via l'ajout d'un vecteur de direction (\emph{steering vector}), de \emph{dropout} ou de bruit gaussien, ainsi que la sensibilité des modèles à l'intensité brute des perturbations, pour compléter les travaux de Fornasiere et al. \cite{fornasiere2026}. Nous étudierons également la sensibilité à l'intensité relative par rapport aux activations nominales, en nous inspirant de Nguyen et al. \cite{nguyen2026}, afin d'évaluer si les différences de détectabilité entre perturbations conceptuelles, aléatoires et bruitées subsistent après cette normalisation. Afin d'obtenir une comparaison plus robuste, toutes les expériences seront menées sur plusieurs architectures différentes, avec une méthodologie rigoureuse mêlant celle de Hahami et al. \cite{hahami2026} à celle décrite par Singh et al. \cite{singh2026}. Enfin, nous étudierons la capacité des modèles à localiser ou à détecter de multiples perturbations, thème peu exploré dans notre état de l'art.

\clearpage
% La largeur et l'orientation sont adaptées au tableau à cinq colonnes.
\newgeometry{left=1.25cm,right=1.25cm,top=1.25cm,bottom=1.25cm}
\begin{landscape}
\section*{Tableau des ressources bibliographiques}
\scriptsize
\setlength{\tabcolsep}{3pt}
\renewcommand{\arraystretch}{1.05}
\begin{longtable}{@{}L{4.5cm} L{7.2cm} L{6.3cm} L{6.3cm}@{}}
\textbf{Travail} & \textbf{Résultat utile} & \textbf{Apport au projet} & \textbf{Limite de l'article}\\
\midrule
\endfirsthead
\textbf{Travail} & \textbf{Résultat utile} & \textbf{Apport au projet} & \textbf{Limite de l'article}\\
\midrule
\endhead
\bottomrule
\endfoot

Hahami et al. \cite{hahami2026}
  \newline \href{https://arxiv.org/abs/2512.12411}{arXiv:2512.12411}
& Les expériences avec choix binaire donné à un modèle sont biaisées dans le cas des injections conceptuelles. Llama 3.1 8B peut détecter et localiser un \emph{steering} conceptuel sur des couches en début d'architecture.
& Méthodologie des expériences avec 2 alternatives et un choix forcé.
& Pas de normalisation de l'amplitude de la perturbation par rapport aux variations usuelles des activations.
\\
\midrule

Macar et al. \cite{macar2026}
  \newline \href{https://arxiv.org/abs/2603.21396}{arXiv:2603.21396}
& L'article explore d'où provient le mécanisme d'introspection dans l'architecture du modèle, comment elle est liée aux vecteurs de concepts injectés, et comment augmenter cette capacité avec trois méthodes : fine-tuning, abliteration et \emph{trained bias vector}.
& Méthodes pour identifier quelles \emph{features} internes du modèle participent au mécanisme d'introspection, et pour augmenter la capacité d'introspection.
& Rien ne dit que les « features » trouvées dans l'article sont vraiment dédiées à l'introspection dans le modèle. De plus, les méthodologies de détection et d'identification sont remises en question dans Hahami et al.
\\
\midrule

Lindsey \cite{lindsey2026}
  \newline (Anthropic), arXiv:2601.01828
& Article fondateur sur le sujet de l'introspection des LLM : définition rigoureuse du concept, cadre et expériences permettant d'évaluer le phénomène.
& Quatre premières expériences, dont trois utilisant du \emph{steering} conceptuel, afin d'évaluer les capacités d'introspection des modèles.
& Malgré quelques résultats notables et surprenants, la grande majorité des situations n'ont mené à aucune observation. Seuls les modèles les plus avancés ont présenté des signaux intéressants.
\\
\midrule

Fornasiere et al. \cite{fornasiere2026}
  \newline (LawZero), arXiv:2604.17465
& Les modèles de langage détectent la présence d'injections de \emph{dropout} ou de bruit gaussien, mais ont des difficultés à classifier ces deux types d'injections.
& Expérience de contrôle détaillée permettant d'évaluer des intensités de perturbations significatives à étudier.
& Résultat décevant sur la classification. Grande disparité d'un modèle à l'autre sur l'ensemble des expériences.
\\
\midrule

Nguyen et al. \cite{nguyen2026}
  \newline (Rice University), arXiv:2605.01167
& Le \emph{steering} produit des perturbations mesurables sur les directions non ciblées, fortement corrélées à la dégradation des performances ($r < -0{,}9$). Un même coefficient ou une même norme de \emph{steering} ne correspond donc pas nécessairement à la même intensité réelle de perturbation.
& Fournit une manière de quantifier et de contrôler l'impact réel d'une injection de \emph{steering}.
& La validation porte surtout sur du \emph{steering} \emph{harmful/harmless}; la généralisation à nos concepts d'introspection reste à vérifier.
\\
\midrule

Mishra et al. \cite{mishra2026}
  \newline (SciForDL), arXiv:2604.09839
& Le \emph{steering} d'activation change la distribution du modèle d'une manière qui n'est pas faisable par des prompts textuels.
& L'article tendrait à dire que ce que l'on mesure n'est pas un phénomène d'introspection, mais la pure détection d'une anomalie.
& L'article adopte une posture très théorique dans sa démonstration. La limite proposée est de ne pas avoir étudié les espaces des modèles quantifiés.
\\
\midrule

Singh et al. \cite{singh2026}
  \newline (NYU), arXiv:2605.26242
& Deux conditions nécessaires sont décrites pour que les résultats d'une expérience proviennent de l'introspection au sens fort et non de calculs simples internes dépendant des entrées : \emph{privileged access} et \emph{second-order computation}.
& Dans les expériences, les entrées ne doivent contenir ni indices ni anomalies. Il faut aussi vérifier si le LLM se représente réellement son état interne, dans un cadre méthodologique permettant d'analyser les causalités.
& Il manque un test expérimental définitif pour montrer l'introspection, et donc probablement pour identifier un mécanisme interne au LLM qui représente et surveille causalement son propre état.
\\
\midrule

Kowalski et al. \cite{kowalski2026}
  \newline (UK AI Security Institute), arXiv:2608.21664
& Les performances sur le benchmark dépendent énormément du modèle. Les auteurs montrent que les LLM sont capables de changer leurs activations internes pour échapper à certains contrôles et ainsi avoir un comportement malveillant ou non aligné même dans leurs poids.
& Donne une méthode, des métriques et un benchmark pour évaluer les capacités d'introspection du modèle.
& Il n'est pas clair que l'auto-contrôle des activations soit une capacité présente sans fine-tuning, une idée que l'on retrouve dans le papier de Macar et al.
\\
\midrule

Ferrara \cite{ferrara2026}
  \newline (University of Southern California), arXiv:2608.20569
& L'information sur la perturbation est bien présente dans les activations, mais n'est pas accessible au rapport verbal spontané du modèle. Cela suggère que l'échec vient du chemin de lecture ou de restitution, et non de l'absence du signal interne.
& Les auteurs montrent qu'il faut comparer la perturbation conceptuelle à une perturbation aléatoire appariée sur son effet aval, et pas seulement sur sa norme.
& Le résultat négatif porte sur les modèles, sites, doses et protocoles testés, pas sur l'introspection des LLM en général.
\\
\end{longtable}
\end{landscape}
\restoregeometry

\clearpage
\begin{thebibliography}{99}
\bibitem{hahami2026}
E. Hahami, I. Sinha, L. Jain, J. Kaplan, and J. Hahami, ``Detecting the Disturbance: A Nuanced View of Introspective Abilities in LLMs,'' Mar. 1, 2026, arXiv:2512.12411. \href{https://doi.org/10.48550/arXiv.2512.12411}{doi:10.48550/arXiv.2512.12411}.

\bibitem{macar2026}
U. Macar, L. Yang, A. Wang, P. Wallich, E. Ameisen, and J. Lindsey, ``Mechanisms of Introspective Awareness,'' Jun. 10, 2026, arXiv:2603.21396. \href{https://doi.org/10.48550/arXiv.2603.21396}{doi:10.48550/arXiv.2603.21396}.

\bibitem{fornasiere2026}
D. Fornasiere, M. Bronzi, S. Kitts, A. Palmas, Y. Bengio, and O. Richardson, ``Language models recognize dropout and Gaussian noise applied to their activations,'' Apr. 30, 2026, arXiv:2604.17465. \href{https://doi.org/10.48550/arXiv.2604.17465}{doi:10.48550/arXiv.2604.17465}.

\bibitem{lindsey2026}
J. Lindsey, ``Emergent Introspective Awareness in Large Language Models,'' Jan. 5, 2026, arXiv:2601.01828. \href{https://doi.org/10.48550/arXiv.2601.01828}{doi:10.48550/arXiv.2601.01828}.

\bibitem{singh2026}
S. Singh, T. Linzen, and S. Ravfogel, ``Can LLMs Introspect? A Reality Check,'' Aug. 21, 2026, arXiv:2605.26242. \href{https://doi.org/10.48550/arXiv.2605.26242}{doi:10.48550/arXiv.2605.26242}.

\bibitem{mishra2026}
A. Mishra, D. Khashabi, and A. Liu, ``Steered LLM Activations are Non-Surjective,'' May 7, 2026, arXiv:2604.09839. \href{https://doi.org/10.48550/arXiv.2604.09839}{doi:10.48550/arXiv.2604.09839}.

\bibitem{nguyen2026}
T. Nguyen, T. A. Nguyen, S. Alemohammad, and R. G. Baraniuk, ``Minimizing Collateral Damage in Activation Steering,'' May 1, 2026, arXiv:2605.01167. \href{https://doi.org/10.48550/arXiv.2605.01167}{doi:10.48550/arXiv.2605.01167}.

\bibitem{kowalski2026}
M. M. Kowalski, J. Fonseca Rivera, U. Macar, and D. D. Africa, ``Measuring Activation Control in Large Language Models,'' Aug. 21, 2026, arXiv:2608.21664. \href{https://arxiv.org/abs/2608.21664}{arXiv:2608.21664}.

\bibitem{ferrara2026}
E. Ferrara, ``Open-Weight Masked Introspection: Measuring What Language Models Can Report About Their Own Computation,'' Aug. 20, 2026, arXiv:2608.20569. \href{https://arxiv.org/abs/2608.20569}{arXiv:2608.20569}.
\end{thebibliography}

\end{document}
